\chapter{État de l'art}

Ce chapitre présente les travaux sur lesquels s'appuie la construction d'un graphe de connaissances clinique à partir de dossiers médicaux scannés. Il suit les étapes d'une telle chaîne de traitement. La première section retrace l'évolution des méthodes d'extraction de texte, des systèmes de reconnaissance optique de caractères aux modèles multimodaux. La deuxième section porte sur l'annotation clinique, qui fournit le corpus de référence indispensable à l'évaluation. La troisième section présente les méthodes d'extraction d'entités et de relations, des approches classiques aux grands modèles de langage et aux modèles multimodaux, et les compare. La quatrième section est consacrée aux hallucinations des modèles multimodaux lors de l'extraction d'entités et de relations médicales, ainsi qu'aux méthodes proposées pour les détecter et les réduire.

\section{Extraction de texte médical : des OCR traditionnels aux modèles multimodaux}

L'exploitation automatique des dossiers médicaux nécessite généralement une première étape de conversion des documents en texte numérique. En effet, une partie importante de l'information clinique reste présente dans des documents non structurés ou semi-structurés tels que les comptes rendus médicaux, les prescriptions, les résultats biologiques, les formulaires hospitaliers ou les lettres de sortie. Lorsque ces documents sont disponibles sous forme scannée ou sous forme d'image, leur traitement automatique nécessite une étape d'extraction documentaire permettant de transformer le contenu visuel en texte exploitable.

Historiquement, cette tâche a été réalisée à l'aide de systèmes de reconnaissance optique de caractères, appelés \textit{Optical Character Recognition} (OCR). L'évolution récente des méthodes d'intelligence artificielle a progressivement déplacé l'intérêt scientifique vers des approches plus avancées, notamment les systèmes de \textit{Document AI} et les grands modèles de langage multimodaux. Cette évolution traduit le passage d'une simple reconnaissance de caractères à une compréhension plus globale du document.

\subsection{Les systèmes OCR traditionnels}

Les systèmes OCR traditionnels constituent la première génération d'outils utilisés pour convertir automatiquement une image contenant du texte en contenu numérique. Parmi les moteurs les plus utilisés figure Tesseract, un moteur OCR open source décrit par Smith~\cite{smith2007tesseract}. Il enchaîne une analyse de la mise en page, qui repère les composantes connexes de l'image puis les lignes et les mots, et une reconnaissance en deux passes reposant sur un classifieur adaptatif : les mots reconnus avec une bonne confiance lors de la première passe servent à mieux reconnaître les autres lors de la seconde.

Ces systèmes ont montré leur efficacité sur des documents imprimés de bonne qualité. Ils restent toutefois fortement dépendants de la qualité visuelle du document. Lorsque l'image est bruitée, inclinée, floue, faiblement contrastée ou structurée sous forme de tableaux complexes, les performances peuvent diminuer considérablement. Dans le contexte médical, cette limite est d'autant plus gênante que les documents hospitaliers peuvent contenir des tampons, des annotations manuscrites, des abréviations, des tableaux biologiques ou des photocopies de faible qualité.

Les erreurs produites par l'OCR peuvent ensuite se propager aux étapes suivantes du pipeline de traitement, notamment l'extraction d'entités médicales, l'extraction de relations et la construction de graphes de connaissances. La qualité de l'extraction textuelle influence donc directement la qualité des informations médicales structurées obtenues par les systèmes automatiques.

\subsection{OCR neuronaux et Document AI}

Afin de dépasser les limites des OCR classiques, plusieurs approches fondées sur l'apprentissage profond ont été proposées. Ces méthodes ne se limitent plus à la reconnaissance caractère par caractère, mais cherchent à exploiter simultanément le texte, la structure visuelle et la mise en page du document.

LayoutLM, proposé par Xu et al.~\cite{xu2020layoutlm}, constitue une avancée importante dans ce domaine. Le modèle combine le contenu textuel, la position spatiale des mots et les caractéristiques visuelles du document afin d'améliorer les tâches de compréhension documentaire. Cette approche convient bien aux documents structurés, car la position d'une information dans la page peut modifier son interprétation : une valeur n'a pas le même sens selon la colonne d'un tableau dans laquelle elle se trouve.

TrOCR, proposé par Li et al.~\cite{li2023trocr}, adopte une architecture Transformer de bout en bout pour la reconnaissance de texte, en associant un encodeur d'images et un décodeur de texte tous deux pré-entraînés. Contrairement aux OCR traditionnels, cette approche ne dépend pas d'une segmentation explicite des caractères : le modèle apprend directement la correspondance entre l'image d'une ligne de texte et sa transcription. TrOCR opère toutefois sur des lignes déjà détectées et ne traite pas à lui seul la page complète.

Donut, proposé par Kim et al.~\cite{kim2022donut}, introduit une approche dite \textit{OCR-free}. Le modèle traite directement l'image complète du document et génère une sortie structurée sans passer par une étape OCR intermédiaire. Cette approche réduit les erreurs accumulées dans les pipelines classiques et marque une transition importante vers la compréhension globale des documents.

Malgré ces avancées, les systèmes de \textit{Document AI} restent généralement orientés vers la transcription ou la compréhension documentaire générale. Ils ne garantissent pas automatiquement la validité clinique des informations extraites et ne traitent pas explicitement le risque de générer des informations absentes du document source.

\subsection{Les modèles multimodaux comme alternative à l'OCR}

L'apparition des grands modèles de langage multimodaux (\textit{Multimodal Large Language Models}, MLLMs) a profondément renouvelé les approches d'extraction documentaire. Ces modèles combinent vision artificielle, traitement du langage naturel et raisonnement contextuel. Contrairement aux OCR classiques, ils peuvent analyser directement une image de document et produire une transcription ou une sortie structurée en tenant compte du contexte global. Parmi les modèles les plus étudiés figurent Gemini, GPT-4o, Qwen-VL, Phi-4 Multimodal et InternVL.

Neveditsin et al.~\cite{neveditsin2026compact} comparent huit systèmes sur des rapports d'échographie obstétricale photographiés, rédigés dans un anglais médical riche en abréviations : des OCR classiques et neuronaux (Tesseract, PaddleOCR, docTR, Surya, GOT-OCR 2.0) et trois MLLMs compacts (Qwen-2.5 VL, Phi-4 Multimodal et InternVL-3.5). Les MLLMs obtiennent les meilleurs taux d'erreur par caractère (\textit{Character Error Rate}, CER) et par mot (\textit{Word Error Rate}, WER). Qwen-2.5 VL obtient ainsi un CER de 0,031 et un WER de 0,078, contre respectivement 0,189 et 0,276 pour Tesseract. La précision sur les valeurs numériques atteint 95,0\,\% pour Qwen-2.5 VL, contre 75,6\,\% pour Tesseract, ce qui est crucial dans le contexte médical, où une erreur sur une valeur biologique peut avoir des conséquences cliniques graves. Ce gain a toutefois un coût : le traitement d'une image prend près d'une minute avec Qwen-2.5 VL, contre moins d'une seconde avec Tesseract.

Inoue~\cite{inoue2025context} évalue GPT-4o et Gemini 2.0 Flash face à l'OCR industriel d'Azure sur la reconnaissance de caractères isolés, des kanji japonais présentés sans aucun contexte, afin de mesurer la capacité de lecture visuelle indépendamment des connaissances linguistiques du modèle. Les MLLMs égalent l'OCR classique à haute résolution (environ 300 ppi), mais leurs performances chutent fortement en dessous de 150 ppi, où ils passent sous l'OCR de référence. Ce résultat souligne la sensibilité des MLLMs à la qualité visuelle des images, et donc à la résolution choisie lors de la numérisation.

Greif et al.~\cite{greif2025multimodal} étudient l'utilisation des MLLMs pour l'OCR, la correction des sorties d'un OCR et la reconnaissance d'entités nommées, sur des documents historiques (annuaires urbains allemands des XVIII\textsuperscript{e} et XIX\textsuperscript{e} siècles). Gemini 2.0 Flash y obtient le meilleur taux d'erreur par caractère, et la correction multimodale de la transcription, qui fournit au modèle à la fois l'image et le texte issu de l'OCR, ramène ce taux sous 1\,\%, sans prétraitement ni affinage. Les auteurs montrent également que ces modèles peuvent extraire directement des entités nommées structurées à partir de ces documents.

Mistral Small 4~\cite{mistralsmall4}, publié par Mistral AI en mars 2026 sous licence Apache 2.0, est le premier modèle de cet éditeur à réunir dans un seul modèle des capacités de vision, de raisonnement et de génération de code, auparavant réparties entre plusieurs modèles spécialisés. Il accepte en entrée du texte et des images, ce qui en fait un MLLM, et repose sur une architecture à mélange d'experts : sur 119 milliards de paramètres au total, environ 6 milliards sont activés pour chaque token, avec une fenêtre de contexte de 256\,000 tokens. Les résultats publiés par l'éditeur portent sur le raisonnement et la programmation, où le modèle égale ou dépasse des modèles ouverts plus volumineux tout en produisant des réponses plus courtes. En revanche, aucune évaluation publiée ne porte à ce jour sur la reconnaissance de documents ou sur des textes médicaux. Ses capacités multimodales, sa licence ouverte et la possibilité de le déployer localement en font néanmoins un candidat pertinent pour l'extraction d'informations à partir de documents médicaux scannés, sous réserve d'une évaluation spécifique.

Ces travaux montrent que les modèles multimodaux constituent des alternatives crédibles aux OCR traditionnels. Toutefois, ils introduisent un risque nouveau : les hallucinations. Contrairement aux OCR classiques, qui se limitent généralement à restituer le texte observé, les MLLMs peuvent produire des informations plausibles mais absentes du document source. Dans le domaine médical, ce risque est particulièrement critique, car l'ajout d'une valeur biologique, d'un traitement ou d'un diagnostic non présent dans le document peut compromettre la fiabilité du pipeline clinique.

\textcolor{red}{Il faut évoquer l'aspect RGPD et anonymisation des données et donc nécessité de tourner en local }

\subsection{Analyse critique des approches d'extraction documentaire}

L'analyse de la littérature montre une évolution progressive des approches d'extraction documentaire. Les OCR traditionnels sont rapides et largement utilisés, mais ils restent fragiles face aux documents bruités. Les modèles de \textit{Document AI} améliorent la prise en compte de la structure visuelle et de la mise en page, mais ils restent principalement orientés vers la compréhension documentaire générale. Les MLLMs offrent une meilleure capacité d'interprétation contextuelle, mais ils sont plus coûteux, sensibles à la résolution des images et exposés à un risque d'hallucination qui nécessite une validation externe.

Par ailleurs, les comparaisons existantes portent sur des corpus restreints et hétérogènes (rapports d'échographie, caractères isolés, documents historiques), et elles évaluent principalement la qualité de transcription, mesurée par le CER et le WER. Ces mesures ne distinguent pas une erreur de lecture d'une information inventée par le modèle. Peu d'études analysent en outre l'impact de ces choix technologiques sur l'ensemble d'un pipeline clinique incluant l'annotation médicale, l'extraction d'entités, l'extraction de relations et la construction d'un graphe de connaissances. Cette limite pèse fortement dans le contexte du sepsis, où la fiabilité de chaque information extraite conditionne la qualité du raisonnement clinique ultérieur.

\section{Annotation clinique, guidelines et pré-annotation}

L'extraction de texte ne suffit pas à rendre les documents médicaux directement exploitables par un système d'intelligence artificielle. Une fois le texte extrait, il est nécessaire d'identifier les informations cliniques pertinentes, de les structurer et de les valider. Cette étape repose sur l'annotation clinique, qui constitue la base de la construction des corpus de référence, de l'entraînement des modèles et de leur évaluation.

\subsection{Annotation clinique et corpus de référence}

L'annotation clinique consiste à identifier dans un document médical les éléments pertinents pour une tâche donnée. Ces éléments peuvent correspondre à des symptômes, des maladies, des infections, des médicaments, des résultats biologiques, des procédures ou encore des relations entre concepts. L'objectif est de transformer un texte libre en données structurées exploitables pour l'entraînement, l'évaluation ou la validation des modèles.

Un \textit{gold standard} correspond à un corpus annoté manuellement et validé par des experts du domaine. Il constitue la vérité de terrain (\textit{ground truth}) utilisée comme référence pour entraîner et évaluer les systèmes automatiques d'extraction d'information. Deleger et al.~\cite{deleger2012building} décrivent la construction de trois corpus de référence de ce type, portant sur des notes cliniques, des annonces d'essais cliniques et des notices de médicaments, et détaillent le processus suivi pour en garantir la fiabilité : guidelines écrites, formation des annotateurs, mesure régulière de l'accord et sessions de consensus.

Dans un contexte clinique, la qualité du corpus de référence est déterminante, car les erreurs d'annotation faussent directement l'estimation des performances. Un système peut sembler performant si la référence utilisée est imprécise, ou au contraire être injustement pénalisé si les annotations humaines sont incohérentes. La qualité du \textit{gold standard} constitue donc une condition essentielle pour toute évaluation expérimentale rigoureuse.

\subsection{Difficultés de l'annotation médicale}

L'annotation médicale présente plusieurs difficultés spécifiques. Les textes cliniques sont souvent courts, télégraphiques, remplis d'abréviations, de synonymes, de négations et d'expressions implicites. Dans leur revue de 174 publications consacrées à l'extraction d'information à partir des dossiers médicaux électroniques, Meystre et al.~\cite{meystre2008extracting} soulignent les difficultés propres au texte clinique, qui imposent des traitements dédiés, notamment pour la détection du contexte d'une mention (négation, temporalité), et relèvent que les systèmes développés sont encore rarement utilisés en dehors des laboratoires qui les ont conçus.

Par exemple, les expressions \og sepsis \fg{}, \og suspicion de sepsis \fg{} et \og antécédent de sepsis \fg{} ne doivent pas être interprétées de la même manière. La première peut indiquer un diagnostic présent, la deuxième une hypothèse clinique et la troisième une information historique. De même, l'expression \og absence de fièvre \fg{} contient le terme \og fièvre \fg{}, mais ne doit pas être interprétée comme un symptôme présent.

Ces difficultés rendent nécessaire la définition de règles explicites. En l'absence de règles communes, deux annotateurs peuvent produire des annotations différentes pour un même document. Cette variabilité réduit la cohérence du corpus et complique l'évaluation des systèmes automatiques. Elle souligne ainsi la nécessité de disposer de règles d'annotation précises et partagées par l'ensemble des annotateurs.

\subsection{Rôle des guidelines d'annotation}

Les guidelines d'annotation sont des documents méthodologiques décrivant les règles que les annotateurs doivent suivre. Elles définissent les catégories à annoter, les critères d'inclusion et d'exclusion, les cas ambigus, les relations autorisées, ainsi que des exemples positifs et négatifs.

Pustejovsky et Stubbs~\cite{pustejovsky2012natural} placent les guidelines au cœur de tout projet d'annotation et les inscrivent dans un cycle itératif, appelé MATTER (\textit{Model, Annotate, Train, Test, Evaluate, Revise}) : un modèle de la tâche est défini, puis appliqué par les annotateurs ; les annotations servent à entraîner et à tester un système, dont l'évaluation conduit à réviser le modèle et les guidelines. Cette démarche permet de transformer une tâche en partie subjective en une procédure contrôlée et reproductible. Dans le domaine médical, le rôle des guidelines est encore plus important, car une même expression peut avoir des interprétations différentes selon le contexte clinique.

Pour un corpus consacré au sepsis, les guidelines doivent notamment définir les types d'entités à annoter, par exemple les symptômes, les foyers infectieux, les micro-organismes, les traitements ou les résultats biologiques, ainsi que les relations autorisées entre ces types, par exemple un traitement administré au patient ou un biomarqueur qui témoigne d'une défaillance d'organe. Elles doivent également préciser la gestion de la négation, de l'incertitude et de la temporalité.

Une guideline efficace ne se limite pas à une simple liste de labels. Elle décrit précisément les critères d'annotation, les cas particuliers et les règles de décision afin de garantir une interprétation homogène des documents par l'ensemble des annotateurs. Elle constitue ainsi un véritable document de référence, qui assure la cohérence des annotations et réduit les désaccords entre experts.

\subsection{Construction itérative des guidelines}

Les guidelines ne sont généralement pas parfaites dès leur première version. Elles doivent être améliorées progressivement à partir des erreurs observées lors des premières sessions d'annotation. Cette démarche itérative est essentielle dans le domaine médical, où certaines ambiguïtés ne deviennent visibles qu'au moment de l'annotation réelle des documents.

Becker et al.~\cite{becker2025cardio} étendent ainsi les guidelines du corpus clinique allemand CARDIO:DE à de nouvelles catégories (diagnostics, médicaments, constatations médicales, thérapies) selon un processus d'annotation itératif. Les auteurs rapportent un accord inter-annotateurs supérieur à 0,9 pour le schéma révisé, et pouvant atteindre 0,98 sur certains sous-ensembles.

Yan et al.~\cite{yan2023terminology} proposent une méthodologie pertinente pour le sepsis et les événements indésirables. Les auteurs construisent une terminologie d'annotation, AAENOTE (\textit{Annotated Adverse Event NOte TErminology}), à partir de discussions avec des cliniciens, et définissent des relations telles que \og une personne présente un signe \fg{}, \og une procédure utilise un dispositif \fg{} ou \og un signe est situé à proximité d'une localisation \fg{}. Leur guide d'annotation a été affiné au fil de quatre itérations, au cours desquelles huit annotateurs ont annoté 100 documents synthétiques décrivant des événements indésirables. La terminologie obtenue comprend 149 classes, 5 propriétés objet, 27 propriétés de données et 4\,470 individus.

Au-delà de ces chiffres, ce travail montre l'intérêt d'une collaboration étroite entre experts cliniques et ingénieurs des connaissances, les ambiguïtés relevées à chaque itération servant à réviser le guide avant la suivante.

\subsection{Pré-annotation automatique}

La création d'un corpus de référence médical est une tâche coûteuse, car elle nécessite du temps et une expertise médicale spécialisée. Pour réduire cette charge, plusieurs travaux ont étudié la pré-annotation automatique. Le principe consiste à utiliser un système automatique pour proposer une première version des annotations, qui est ensuite corrigée ou validée par les experts.

Lingren et al.~\cite{lingren2014preannotation} évaluent l'impact de la pré-annotation sur la vitesse d'annotation et sur le risque de biais, dans le cadre de la reconnaissance d'entités cliniques dans des annonces d'essais cliniques. La pré-annotation réduit le temps d'annotation de 13,9\,\% à 21,5\,\% selon les entités, sans dégradation significative de l'accord inter-annotateurs, qui reste compris entre 93,4\,\% et 95,5\,\%. Le risque que les annotateurs reprennent sans les corriger les erreurs de la pré-annotation a été explicitement évalué, sans qu'un biais significatif soit observé.

Avec l'émergence des LLMs et des MLLMs, la pré-annotation peut être réalisée par des modèles capables d'identifier des entités et des relations médicales à partir du texte, voire directement à partir d'une image de document. Hors du domaine médical, Kholodna et al.~\cite{kholodna2024llms} montrent, pour la reconnaissance d'entités nommées dans des langues africaines peu dotées, qu'un LLM (GPT-4-Turbo) intégré à une boucle d'apprentissage actif permet d'atteindre des performances proches de l'état de l'art pour un coût d'annotation nettement inférieur, tout en restant en deçà de l'annotation humaine. En contexte médical, l'utilisation d'un LLM comme pré-annotateur appelle toutefois une validation experte systématique, car ces modèles peuvent produire des annotations plausibles mais absentes du document source. Ce risque d'hallucination est étudié en détail dans la suite de ce chapitre.

\subsection{Accord inter-annotateurs et qualité du corpus}

La qualité d'un corpus annoté est généralement évaluée à l'aide de mesures d'accord inter-annotateurs, qui permettent de vérifier si plusieurs experts appliquent les guidelines de manière cohérente. Le coefficient kappa de Cohen (\textit{Cohen's kappa}) est la mesure la plus répandue pour les annotations catégorielles. Hripcsak et Rothschild~\cite{hripcsak2005agreement} montrent cependant qu'il n'est pas calculable pour les tâches d'extraction d'entités, car le nombre de cas négatifs, c'est-à-dire de segments de texte qui ne constituent pas une entité, n'y est pas défini. La F-mesure calculée entre deux annotateurs, en prenant l'un comme référence, est alors la mesure appropriée, et le kappa tend vers elle lorsque les négatifs sont très nombreux.

Deleger et al.~\cite{deleger2012building} évaluent ainsi l'accord inter-annotateurs à l'aide de la F-mesure. Grâce à un suivi régulier de l'accord et à des sessions de consensus, ils obtiennent un accord élevé, compris entre 0,8467 et 0,9176 selon les tâches (désidentification des informations personnelles, extraction des médicaments, des maladies et des signes cliniques), ce qui montre qu'une démarche structurée d'amélioration des guidelines permet d'obtenir un corpus de référence fiable.

James~\cite{james2026consensus} propose une revue des mesures d'accord organisée par type de tâche (étiquetage catégoriel, segmentation, jugement subjectif, notation continue) et rappelle que le choix de la mesure doit être guidé par la nature de la tâche d'annotation, en tenant compte du déséquilibre entre les étiquettes et des données manquantes. Il est donc essentiel de préciser la mesure utilisée ainsi que son mode d'interprétation afin de garantir une évaluation rigoureuse de la qualité du corpus.

Dans le cas de l'annotation médicale appliquée au sepsis, un accord élevé entre annotateurs garantit que les concepts cliniques et les relations sont identifiés de manière homogène avant leur utilisation pour l'entraînement ou l'évaluation des systèmes d'extraction d'information.

\subsection{Analyse critique des travaux sur l'annotation clinique}

L'analyse de la littérature montre que la qualité des systèmes de traitement automatique des langues en domaine médical dépend fortement de la qualité des annotations, des guidelines et des corpus de référence. Les travaux existants démontrent qu'une démarche itérative, associant experts du domaine et spécialistes du traitement automatique des langues, permet d'améliorer progressivement la cohérence des annotations et la fiabilité des corpus de référence. Par ailleurs, les approches de pré-annotation contribuent à réduire l'effort humain tout en maintenant un niveau élevé d'accord inter-annotateurs, sous réserve d'une validation experte.

Plusieurs limites subsistent cependant. La majorité des travaux portent sur des documents nativement numériques ; peu d'études s'intéressent aux documents médicaux scannés, pour lesquels les erreurs d'extraction du texte peuvent se propager jusqu'à l'étape d'annotation. De plus, rares sont les approches qui combinent simultanément extraction multimodale, pré-annotation assistée par un LLM, validation experte, guidelines spécialisées pour le sepsis et construction d'un graphe de connaissances clinique. Ces limites mettent en évidence la nécessité de méthodologies intégrées garantissant la fiabilité de l'ensemble de la chaîne de traitement, depuis l'extraction documentaire jusqu'à la structuration des connaissances médicales.

\section{Extraction d'entités et de relations cliniques : des méthodes classiques aux MLLMs}
\label{sec:extraction_er}

Une fois le texte disponible et le schéma d'annotation défini, l'étape suivante consiste à extraire automatiquement les informations cliniques et à les structurer. Cette tâche se décompose en deux sous-tâches. La reconnaissance d'entités nommées (\textit{Named Entity Recognition}, NER) repère dans le texte les segments qui désignent un concept d'intérêt et leur attribue un type, par exemple \og pneumopathie \fg{} comme diagnostic ou \og noradrénaline \fg{} comme traitement. L'extraction de relations (\textit{Relation Extraction}, RE) identifie ensuite les liens entre ces entités, par exemple un traitement administré au patient ou un micro-organisme isolé dans un prélèvement. Les entités et les relations ainsi obtenues forment les nœuds et les arêtes d'un graphe de connaissances.

Les méthodes proposées pour ces deux tâches ont suivi l'évolution générale du traitement automatique des langues : règles et dictionnaires, apprentissage automatique supervisé, modèles de langage pré-entraînés, puis grands modèles de langage génératifs et, plus récemment, modèles multimodaux capables de travailler directement sur l'image du document. Cette section présente ces familles, puis les compare.

\subsection{Approches à base de règles et de dictionnaires}

Les premiers systèmes d'extraction clinique reposent sur des terminologies médicales et sur des règles écrites par des experts. MetaMap, développé par la National Library of Medicine et présenté par Aronson et Lang~\cite{aronson2010metamap}, associe les segments d'un texte biomédical aux concepts de l'UMLS. Les auteurs retracent une dizaine d'années de développement, au cours desquelles se sont ajoutées la détection des acronymes définis dans le texte, la détection de la négation et la désambiguïsation du sens des mots. cTAKES, proposé par Savova et al.~\cite{savova2010ctakes}, est un système open source construit sur l'architecture UIMA, qui enchaîne des composants entraînés pour le domaine clinique (découpage en phrases, étiquetage morphosyntaxique, analyse en groupes syntaxiques). La reconnaissance des entités y repose sur une recherche dans un dictionnaire, dérivé de l'UMLS et couvrant SNOMED CT et RxNorm, à l'intérieur des groupes nominaux. Le système obtient une F-mesure de 0,715 en correspondance exacte des segments et de 0,824 en correspondance par recouvrement.

Ces approches présentent des qualités durables : elles ne nécessitent pas de corpus annoté, leurs décisions sont explicables et elles relient directement les entités extraites à une terminologie normalisée. Elles restent toutefois limitées par la couverture des dictionnaires, la variabilité de l'écriture clinique (abréviations, fautes, formulations télégraphiques) et le coût de maintenance des règles. Elles traitent en outre surtout les entités : l'extraction de relations y demande des règles supplémentaires, propres à chaque type de relation.

\subsection{Apprentissage automatique supervisé}

L'apprentissage supervisé a ensuite permis d'apprendre les règles à partir de corpus annotés. Pour la reconnaissance d'entités, la tâche est formulée comme un étiquetage de séquence : chaque mot reçoit une étiquette indiquant s'il commence, continue ou non une entité d'un type donné. Les champs aléatoires conditionnels (\textit{Conditional Random Fields}, CRF), introduits par Lafferty et al.~\cite{lafferty2001crf}, sont devenus le modèle de référence pour cette formulation, car ils modélisent les dépendances entre les étiquettes voisines tout en combinant librement des caractéristiques du texte (forme du mot, préfixes, appartenance à un dictionnaire). L'extraction de relations est quant à elle généralement traitée comme une classification de paires d'entités.

Les campagnes d'évaluation partagées ont joué un rôle central dans le domaine clinique. Le défi i2b2/VA 2010~\cite{uzuner2011i2b2} portait sur trois tâches menées sur des comptes rendus de sortie et des notes d'évolution : l'extraction des concepts médicaux (problèmes, examens, traitements), la classification de leur statut d'assertion (présent, absent, possible, hypothétique, conditionnel, concernant une autre personne) et la classification des relations entre concepts. Le système du Conseil national de recherches du Canada, décrit par de Bruijn et al.~\cite{debruijn2011machine}, se classe premier en extraction des concepts (F-mesure de 0,8523) et en classification des assertions (0,9362), et deuxième en classification des relations (0,7313). Il repose sur l'apprentissage supervisé et sur un très grand nombre de caractéristiques, tirées du texte lui-même et de ressources externes (UMLS, cTAKES, MEDLINE) ; les auteurs concluent que la diversité de ces caractéristiques a été déterminante. Ce défi montre également que le statut d'une mention, déjà évoqué à propos de l'annotation, est traité comme une tâche à part entière.

Avec l'apprentissage profond, les caractéristiques conçues à la main ont été remplacées par des représentations apprises. Lample et al.~\cite{lample2016neural} proposent une architecture combinant un réseau récurrent bidirectionnel (BiLSTM) et une couche CRF, qui exploite des représentations des mots construites à partir de leurs caractères et atteint l'état de l'art dans quatre langues sans ressource lexicale externe. Cette architecture a été largement reprise pour la reconnaissance d'entités cliniques.

\subsection{Modèles de langage pré-entraînés}

Les modèles de langage pré-entraînés fondés sur l'architecture Transformer, à commencer par BERT~\cite{devlin2019bert}, ont ensuite renouvelé l'état de l'art : un modèle est pré-entraîné sur de grandes quantités de texte non annoté, puis affiné sur la tâche cible avec un corpus annoté de taille modeste. Des variantes spécialisées ont été pré-entraînées sur des textes du domaine. BioBERT~\cite{lee2020biobert}, entraîné sur des articles biomédicaux, améliore l'état de l'art précédent de 0,62 point de F-mesure en reconnaissance d'entités biomédicales et de 2,80 points en extraction de relations. Alsentzer et al.~\cite{alsentzer2019clinicalbert} poursuivent ce pré-entraînement sur les notes cliniques de MIMIC-III : leur modèle Bio+ClinicalBERT obtient une F-mesure de 87,2 sur l'extraction des concepts du défi i2b2 2010, contre 83,5 pour BERT générique.

Pour le français, DrBERT~\cite{labrak2023drbert} est pré-entraîné sur NACHOS, un corpus de plus d'un milliard de mots collectés sur des sites médicaux francophones, et évalué sur onze tâches biomédicales françaises ; les auteurs comparent également des données publiques issues du web à des données hospitalières privées. CamemBERT-bio~\cite{touchent2024camembertbio} poursuit le pré-entraînement du modèle généraliste CamemBERT sur 413 millions de mots de textes biomédicaux publics et améliore en moyenne de 2,54 points de F-mesure les résultats en reconnaissance d'entités biomédicales françaises. Les ressources annotées en français restent cependant rares : le corpus QUAERO~\cite{neveol2014quaero}, l'une des références, compte environ 100\,000 mots issus de notices de médicaments, de titres d'articles MEDLINE et de brevets, et non de dossiers patients.

Ces modèles affinés constituent aujourd'hui la référence lorsque des données annotées sont disponibles. Lors du défi n2c2 2018~\cite{henry2020n2c2}, portant sur l'extraction des médicaments et des événements indésirables dans 505 comptes rendus de sortie de MIMIC-III, la majorité des systèmes reposaient sur l'apprentissage profond, et les meilleurs atteignaient une F-mesure de 0,9418 pour l'extraction des entités, de 0,9630 pour la classification des relations lorsque les entités de référence sont fournies, et de 0,8905 pour l'extraction de bout en bout. L'écart entre ces deux derniers chiffres illustre un phénomène général : les erreurs commises sur les entités se propagent aux relations.

La revue systématique de Fraile Navarro et al.~\cite{fraile2023review}, portant sur 94 études consacrées à la reconnaissance d'entités et à l'extraction de relations dans les textes cliniques, montre que la majorité des travaux reposent sur l'apprentissage automatique, que les approches à base de règles restent minoritaires, et que seules trois études rapportent une utilisation en dehors d'un cadre expérimental. Les performances obtenues sur les benchmarks ne garantissent donc pas, à elles seules, une utilisation en pratique clinique.

\subsection{Extraction en chaîne et extraction conjointe}

Deux architectures coexistent pour combiner les deux sous-tâches. L'approche en chaîne (\textit{pipeline}) reconnaît d'abord les entités, puis classe les relations entre les paires obtenues ; elle est simple, mais une entité manquée rend impossible toute relation qui l'implique. L'approche conjointe extrait entités et relations au sein d'un même modèle afin que les deux décisions s'informent mutuellement. Miwa et Bansal~\cite{miwa2016endtoend} proposent ainsi un modèle de bout en bout qui superpose un réseau récurrent sur l'arbre de dépendances à un réseau récurrent sur la séquence, et réduisent l'erreur relative de 12,1\,\% sur ACE2005 par rapport au meilleur système à caractéristiques manuelles. SpERT~\cite{eberts2020spert} réalise l'extraction conjointe à partir des segments candidats, à l'aide d'un modèle BERT. La supériorité de l'approche conjointe n'est toutefois pas acquise : Zhong et Chen~\cite{zhong2021frustratingly} montrent qu'une chaîne utilisant deux encodeurs indépendants, l'un pour les entités et l'autre pour les relations, dépasse les modèles conjoints de 1,7 à 2,8 points de F-mesure sur les relations, à encodeur pré-entraîné identique. Le travail LATTE~\cite{zhang2026latte}, présenté plus loin, s'inscrit dans cette lignée en ciblant la stabilité de l'extraction conjointe lorsque plusieurs relations partagent une même entité.

\subsection{Extraction par les grands modèles de langage}

Les grands modèles de langage génératifs abordent l'extraction d'une autre manière : au lieu d'étiqueter chaque mot, ils génèrent directement la liste des entités et des relations, souvent dans un format structuré comme JSON, à partir d'une invite qui décrit la tâche. Xu et al.~\cite{xu2024genie_survey} recensent ces approches génératives d'extraction d'information. Leur principal intérêt est de pouvoir fonctionner sans exemple (\textit{zero-shot}) ou avec quelques exemples (\textit{few-shot}), donc sans corpus d'entraînement, et de s'adapter à un nouveau schéma en modifiant seulement l'invite. Agrawal et al.~\cite{agrawal2022fewshot} ont montré que GPT-3 constitue un bon extracteur d'information clinique dans ces conditions, sans avoir été entraîné sur des données cliniques, à condition de guider le format de sortie et de post-traiter les réponses. Wei et al.~\cite{wei2023chatie} proposent ChatIE, qui décompose l'extraction en un dialogue en plusieurs tours avec ChatGPT : le modèle identifie d'abord les types présents, puis extrait les éléments de chaque type, ce qui lui permet de dépasser certains modèles entièrement supervisés sur plusieurs jeux de données.

Les comparaisons systématiques nuancent toutefois cet enthousiasme. Gutiérrez et al.~\cite{gutierrez2022gpt3bioie} montrent, sur l'ensemble des tâches de reconnaissance d'entités et d'extraction de relations du benchmark biomédical BLURB, que l'apprentissage en contexte de GPT-3 reste nettement inférieur à l'affinage d'un modèle de la taille de BERT, même avec une calibration et une sélection dynamique des exemples. Chen et al.~\cite{chen2025benchmarking} évaluent quatre LLMs (GPT-3.5, GPT-4, LLaMA~2 et PMC-LLaMA) sur douze benchmarks biomédicaux couvrant six applications : les modèles affinés dépassent les LLMs utilisés sans ou avec quelques exemples sur dix des douze jeux de données. L'écart est particulièrement marqué en reconnaissance d'entités : sur NCBI-Disease, la meilleure F-mesure publiée est de 0,9090, contre 0,5988 pour GPT-4. Les auteurs relèvent aussi des réponses incohérentes, manquantes ou hallucinées, ainsi que des erreurs de délimitation des entités.

La qualité de l'invite réduit une partie de cet écart. Hu et al.~\cite{hu2024clinicalner} proposent pour la reconnaissance d'entités cliniques une invite en quatre parties : description de la tâche et du format, extraits des guidelines d'annotation, instructions tirées de l'analyse des erreurs, et exemples annotés. Sur le corpus MTSamples, la F-mesure souple de GPT-4 passe de 0,804 avec une invite de base à 0,861 avec l'invite complète, sans atteindre celle de BioClinicalBERT affiné (0,901) ; sur le corpus VAERS, elle passe de 0,593 à 0,736, contre 0,802. Sainz et al.~\cite{sainz2024gollie} généralisent l'idée avec GoLLIE, un modèle affiné pour suivre des guidelines d'annotation fournies dans l'invite sous forme de définitions de classes commentées. Sur des jeux de données jamais vus, l'ajout des guidelines apporte en moyenne 13 points de F-mesure par rapport au même modèle sans guidelines, et l'étude d'ablation montre que des définitions détaillées sont déterminantes. Ces deux travaux montrent que des guidelines explicites constituent un levier efficace pour l'extraction par LLM. Une autre voie consiste à transférer les capacités d'un grand modèle vers un modèle plus petit et spécialisé : UniversalNER~\cite{zhou2024universalner} distille ainsi des annotations produites par ChatGPT dans un modèle LLaMA, qui dépasse ChatGPT de 7 à 9 points de F-mesure en moyenne sur 43 jeux de données de 9 domaines, dont le biomédical.

Pour l'extraction de relations, Wadhwa et al.~\cite{wadhwa2023revisiting} soulignent que la correspondance exacte, utilisée pour évaluer les systèmes supervisés, pénalise les modèles génératifs, qui formulent souvent une relation correcte avec d'autres mots que la référence. En évaluation humaine, GPT-3 utilisé avec quelques exemples atteint un niveau comparable à celui des modèles entièrement supervisés. L'extraction conjointe d'entités et de relations sous forme de sortie structurée a également été obtenue par affinage : Dagdelen et al.~\cite{dagdelen2024structured} affinent GPT-3 et LLaMA~2 pour produire directement des objets JSON reliant les entités d'articles de science des matériaux, avec quelques centaines d'exemples d'entraînement seulement.

Les LLMs introduisent enfin des erreurs absentes des modèles d'étiquetage. Un modèle d'étiquetage ne peut annoter que des mots présents dans le texte, alors qu'un modèle génératif peut produire une entité reformulée ou inventée. Wang et al.~\cite{wang2025gptner} observent ainsi, avec GPT-NER, que les LLMs ont tendance à étiqueter avec trop d'assurance comme entités des segments qui n'en sont pas, et ajoutent une étape d'auto-vérification pour limiter ce phénomène. Cette particularité est étudiée en détail dans la section consacrée aux hallucinations.

\subsection{Extraction directe à partir de l'image du document}

Les méthodes précédentes supposent que le texte a déjà été extrait du document. Pour des documents scannés, une autre voie consiste à extraire les informations directement à partir de l'image. Dans le domaine de la compréhension documentaire, l'extraction d'informations clés (\textit{Key Information Extraction}, KIE) consiste à retrouver des champs prédéfinis, comme le montant d'un reçu ou le nom d'un formulaire. LayoutLMv3~\cite{huang2022layoutlmv3}, qui pré-entraîne conjointement le texte et l'image du document, a atteint l'état de l'art sur la compréhension de formulaires (FUNSD) et de reçus (CORD). Donut~\cite{kim2022donut}, présenté dans la première section, réalise cette extraction sans OCR intermédiaire. Ces modèles nécessitent toutefois un affinage sur des données annotées du même type de document.

Les MLLMs permettent d'aborder cette tâche sans affinage, à partir d'une invite. OCRBench~\cite{liu2024ocrbench} évalue leurs capacités de lecture sur 1\,000 paires de questions et réponses vérifiées manuellement, issues de 29 jeux de données et réparties en cinq tâches, dont l'extraction d'informations clés. UniKIE-Bench~\cite{ji2026unikie}, qui évalue quinze modèles multimodaux, montre que leurs performances chutent fortement lorsque le schéma des champs varie, lorsque les champs sont rares ou lorsque la mise en page est complexe. Greif et al.~\cite{greif2025multimodal}, cités précédemment, montrent par ailleurs qu'un MLLM peut extraire directement des entités nommées structurées à partir d'images de documents historiques.

Dans le domaine médical, MedRepBench~\cite{shang2025medrepbench} évalue l'extraction structurée à partir d'environ 1\,900 images de comptes rendus d'examens et de résultats biologiques chinois réels, désidentifiés, en mesurant pour chaque résultat le nom, la valeur, l'unité, l'intervalle de référence et le caractère anormal. Les modèles vision-langage appliqués directement à l'image restent en deçà d'une chaîne combinant un OCR de haute qualité et un LLM textuel, utilisée comme borne supérieure ; les auteurs notent cependant que cette chaîne ignore la mise en page et qu'elle est plus lente. Schäfer et al.~\cite{schaefer2025transfusion} combinent, pour des formulaires papier scannés de déclaration de réactions transfusionnelles, un détecteur de cases à cocher et un modèle vision-langage qui associe le contenu à des catégories normalisées, et évaluent leur système contre les références annuelles de 2017 à 2024. À l'inverse, lorsque le texte passe par un OCR, ses erreurs se propagent aux étapes suivantes : Zhang et al.~\cite{zhang2025ocrhindersrag} montrent, sur 8\,561 images de documents de sept domaines, que le bruit introduit par l'OCR, qu'il porte sur le contenu ou sur la mise en forme, dégrade les systèmes qui s'appuient sur ce texte, et qu'aucune des solutions OCR évaluées ne suffit à construire une base de connaissances fiable.

Ces travaux portent cependant sur l'extraction de champs, c'est-à-dire de paires clé--valeur, et non sur l'extraction de relations typées entre entités. À notre connaissance, aucune étude évaluée par les pairs ne compare, sur des documents cliniques scannés, l'extraction d'entités et de relations par un MLLM appliqué à l'image et par une chaîne OCR suivie d'un modèle textuel.

\subsection{Comparaison des approches}

Le Tableau~\ref{tab:comparaison_er} résume les caractéristiques des différentes familles de méthodes.

{\footnotesize
\renewcommand{\arraystretch}{1.35}
\begin{longtable}{|>{\raggedright\arraybackslash}p{2.1cm}|>{\raggedright\arraybackslash}p{2.4cm}|>{\raggedright\arraybackslash}p{2.2cm}|>{\raggedright\arraybackslash}p{3.0cm}|>{\raggedright\arraybackslash}p{3.4cm}|}
\caption{Comparaison des familles de méthodes d'extraction d'entités et de relations.}
\label{tab:comparaison_er}\\
\hline
\textbf{Famille} & \textbf{Exemples} & \textbf{Données annotées nécessaires} & \textbf{Points forts} & \textbf{Limites} \\
\hline
\endfirsthead
\multicolumn{5}{l}{\textit{Tableau~\ref{tab:comparaison_er} (suite)}}\\
\hline
\textbf{Famille} & \textbf{Exemples} & \textbf{Données annotées nécessaires} & \textbf{Points forts} & \textbf{Limites} \\
\hline
\endhead
\hline
\multicolumn{5}{r}{\textit{Suite à la page suivante}}\\
\endfoot
\hline
\endlastfoot

Règles et dictionnaires &
MetaMap \cite{aronson2010metamap}, cTAKES \cite{savova2010ctakes} &
Aucune (terminologies et règles expertes) &
Décisions explicables, normalisation directe vers l'UMLS, pas d'invention possible &
Couverture limitée des dictionnaires, sensibilité aux abréviations et aux fautes, maintenance coûteuse, relations peu couvertes \\
\hline

Apprentissage supervisé classique &
CRF \cite{lafferty2001crf}, BiLSTM-CRF \cite{lample2016neural} &
Corpus annoté de taille importante &
Apprentissage des variations d'écriture, bonnes performances sur le domaine d'entraînement &
Caractéristiques à concevoir (CRF), dépendance au corpus annoté, transfert difficile vers un autre schéma \\
\hline

Modèles pré-entraînés affinés &
BioBERT \cite{lee2020biobert}, ClinicalBERT \cite{alsentzer2019clinicalbert}, DrBERT \cite{labrak2023drbert}, CamemBERT-bio \cite{touchent2024camembertbio}, PURE \cite{zhong2021frustratingly} &
Corpus annoté de taille modeste &
Meilleures performances actuelles sur les benchmarks annotés \cite{chen2025benchmarking, henry2020n2c2} &
Nouvel affinage à chaque changement de schéma, peu de corpus cliniques annotés en français, texte préalablement extrait \\
\hline

LLMs génératifs (sans exemple ou avec quelques exemples) &
GPT-3 \cite{agrawal2022fewshot}, ChatIE \cite{wei2023chatie}, invites guidées par les guidelines \cite{hu2024clinicalner} &
Aucune ou quelques exemples &
Pas d'entraînement, adaptation au schéma par l'invite, gain notable apporté par les guidelines \cite{hu2024clinicalner, sainz2024gollie} &
En deçà des modèles affinés \cite{gutierrez2022gpt3bioie, chen2025benchmarking}, entités reformulées ou inventées, évaluation par correspondance exacte inadaptée \cite{wadhwa2023revisiting} \\
\hline

LLMs affinés ou distillés &
GoLLIE \cite{sainz2024gollie}, UniversalNER \cite{zhou2024universalner}, Dagdelen et al. \cite{dagdelen2024structured} &
Quelques centaines d'exemples ou données produites par un LLM &
Sortie structurée directe (JSON), extraction conjointe, bonne généralisation &
Coût d'affinage, accès aux poids nécessaire \\
\hline

MLLMs appliqués à l'image &
OCRBench \cite{liu2024ocrbench}, UniKIE-Bench \cite{ji2026unikie}, MedRepBench \cite{shang2025medrepbench} &
Aucune (utilisation par invite) &
Pas d'OCR intermédiaire, prise en compte de la mise en page, une seule étape &
Sensibilité au schéma, à la mise en page et à la qualité de l'image, résultats encore en deçà d'un OCR de qualité suivi d'un LLM \cite{shang2025medrepbench}, évaluation limitée aux champs clé--valeur \\
\end{longtable}
}

\subsection{Analyse critique}

Les méthodes supervisées, et en particulier les modèles pré-entraînés affinés, restent les plus performantes lorsqu'un corpus annoté de taille suffisante est disponible : les études comparatives les plus larges montrent qu'elles dépassent encore les LLMs génératifs utilisés sans affinage, surtout en reconnaissance d'entités~\cite{gutierrez2022gpt3bioie, chen2025benchmarking}. Cette supériorité suppose cependant des conditions rarement réunies en pratique clinique : un corpus annoté dans la langue et le domaine visés, un schéma stable et un texte déjà extrait. Or les corpus cliniques annotés en français sont rares, la plupart des ressources disponibles portant sur des textes biomédicaux publics plutôt que sur des dossiers patients~\cite{neveol2014quaero, touchent2024camembertbio}, et un schéma propre à une pathologie comme le sepsis évolue au fil des itérations des guidelines.

Les LLMs et les MLLMs répondent à ces contraintes en déplaçant l'effort de l'annotation d'un corpus d'entraînement vers la rédaction de l'invite et des guidelines, dont l'apport a été mesuré~\cite{hu2024clinicalner, sainz2024gollie}, et, pour les MLLMs, en évitant la propagation des erreurs d'OCR~\cite{zhang2025ocrhindersrag}. Ils introduisent en contrepartie deux difficultés. La première est l'évaluation : la correspondance exacte pénalise des réponses correctes mais formulées autrement~\cite{wadhwa2023revisiting}, ce qui impose des critères d'appariement plus souples. La seconde est la fidélité : contrairement à un modèle d'étiquetage, qui ne peut sélectionner que des segments présents dans le texte, un modèle génératif peut produire une entité, une valeur ou une relation absente du document~\cite{wang2025gptner, chen2025benchmarking}. Enfin, l'extraction directe d'entités et de relations à partir d'images de documents cliniques n'a pas encore fait l'objet d'une évaluation comparative : les benchmarks multimodaux existants portent sur des champs clé--valeur~\cite{liu2024ocrbench, ji2026unikie, shang2025medrepbench}. Le risque d'hallucination propre à ces modèles, qui conditionne leur utilisation pour construire un graphe de connaissances clinique, fait l'objet de la section suivante.

\section{Hallucination des modèles multimodaux lors de l'extraction d'entités et de relations médicales}

Comme évoqué lors de la présentation des modèles multimodaux, l'utilisation des MLLMs dans les pipelines d'extraction d'information clinique introduit un risque spécifique : celui de générer des entités, des relations ou des valeurs plausibles mais absentes du document source. Ce phénomène, appelé hallucination, découle du fait que ces modèles interprètent le contenu à partir de leurs connaissances apprises, au lieu de se limiter à restituer le texte observé.

Cette section présente les principaux travaux consacrés à la compréhension, à la détection et à la réduction des hallucinations, en s'intéressant surtout à l'extraction structurée d'entités et de relations, où chaque information extraite doit être fidèle au document source pour pouvoir alimenter un graphe de connaissances.

\subsection{Définitions et taxonomies générales de l'hallucination}

La littérature distingue généralement deux grandes catégories d'hallucinations. Ji et al.~\cite{ji2023survey}, dans leur revue consacrée à la génération de texte, opposent les hallucinations intrinsèques, qui correspondent à des informations contredisant directement le contenu du document source, aux hallucinations extrinsèques, qui désignent des informations ne pouvant pas être vérifiées à partir du document, sans pour autant être nécessairement fausses. Huang et al.~\cite{huang2023survey} prolongent cette vision dans le cas des LLMs en distinguant les hallucinations de factualité, qui contredisent les connaissances du monde, des hallucinations de fidélité, qui s'écartent des instructions de l'utilisateur ou du contexte fourni, ou qui sont incohérentes avec elles-mêmes. Ils identifient également les différentes sources d'hallucination tout au long du cycle de vie d'un modèle, réparties en trois étapes : les données, l'entraînement (pré-entraînement, ajustement supervisé et apprentissage par renforcement à partir de retours humains) et l'inférence, et distinguent les méthodes visant à détecter les hallucinations de celles destinées à les atténuer. Dans le cadre de l'extraction d'information, c'est la fidélité qui importe : une information médicalement exacte mais absente du dossier du patient reste une erreur.

Anh-Hoang et al.~\cite{anhhoang2025survey} s'intéressent quant à eux à l'origine des hallucinations au moment de l'utilisation du modèle : ils proposent un cadre d'attribution qui distingue les hallucinations dues à la formulation de l'invite (\textit{Prompt Sensitivity}) de celles dues au comportement intrinsèque du modèle (\textit{Model Variability}). Leurs expériences montrent que des stratégies d'invite structurées, comme le raisonnement pas à pas (\textit{chain-of-thought}), réduisent nettement les hallucinations lorsque celles-ci dépendent de l'invite, mais que certaines limites propres au modèle persistent. Pour une tâche d'extraction, cela signifie que la conception de l'invite est un levier utile mais insuffisant.

Dans le cas des modèles multimodaux, Bai et al.~\cite{bai2025survey} proposent une taxonomie plus spécifique adaptée aux interactions entre les modalités visuelle et textuelle. Leur étude recense les principales causes d'hallucination, liées aux données, au modèle, à l'entraînement ou à l'inférence, les benchmarks utilisés pour les évaluer ainsi que les différentes stratégies de mitigation. Les auteurs distinguent notamment les hallucinations portant sur la catégorie d'un objet, c'est-à-dire sur son existence même, sur ses attributs (par exemple la couleur, la forme ou le nombre) ou encore sur les relations entre plusieurs objets, comme leurs interactions ou leurs positions relatives. Cette classification est particulièrement pertinente pour l'extraction d'informations à partir de documents médicaux numérisés, où une erreur sur une entité, une valeur clinique ou une relation peut se propager aux étapes suivantes du pipeline, notamment lors de la construction d'un graphe de connaissances.

Bien que ces taxonomies constituent une référence pour l'étude des hallucinations, elles ont principalement été élaborées pour des tâches de génération libre, telles que la description d'images, la réponse à des questions ou le résumé automatique. Elles prennent donc peu en compte les spécificités des tâches d'extraction structurée, dans lesquelles les informations produites doivent respecter un schéma prédéfini d'entités et de relations. Cette différence est importante, car les erreurs observées dans ces deux contextes ne sont pas de même nature. En extraction structurée, une hallucination ne se limite pas à l'invention d'une information absente du document ; elle peut également se traduire par une incohérence structurelle, par exemple lorsqu'une relation est créée entre deux entités incompatibles avec les contraintes du schéma ou lorsqu'une entité est associée à un type de relation qui ne lui est pas autorisé. Ce type d'erreur est critique dans la construction d'un graphe de connaissances, où la validité des relations dépend directement de la cohérence entre les types d'entités et les relations qui les relient.

\subsection{Hallucination dans l'extraction d'entités et de relations hors domaine médical}

Les travaux consacrés à la reconnaissance d'entités nommées (NER) et à l'extraction de relations (RE) à l'aide des LLMs évaluent généralement les performances à l'aide de métriques globales, telles que la précision, le rappel et le F1-score. En revanche, ils s'intéressent rarement à la nature des erreurs produites par les modèles et les analysent encore moins sous l'angle des hallucinations.

Tran et al.~\cite{tran2025aspworkflow} proposent un pipeline combinant un LLM et l'\textit{Answer Set Programming} (ASP), un paradigme de programmation déclarative fondé sur la logique, dans lequel les connaissances sont représentées sous forme de faits et de règles permettant de vérifier automatiquement le respect de contraintes prédéfinies. Dans cette approche, le LLM réalise l'extraction initiale des entités et des relations, tandis que le programme ASP applique un raisonnement symbolique afin de contrôler la cohérence des informations extraites, par exemple en vérifiant la compatibilité entre les types d'entités et les relations autorisées. Cette étape permet d'éliminer certaines incohérences structurelles avant la production du résultat final, et des connaissances du domaine peuvent être ajoutées sous forme de règles sans modifier le programme central. L'approche est évaluée sur trois benchmarks d'extraction conjointe d'entités et de relations, avec peu de données annotées. Bien que ces erreurs ne soient pas explicitement qualifiées d'hallucinations par les auteurs, cette approche montre qu'une validation fondée sur des contraintes logiques peut renforcer la fiabilité des résultats produits par le modèle.

Sur le plan méthodologique, plusieurs approches de détection proposées dans la littérature peuvent être adaptées à la tâche d'extraction d'information. Certaines d'entre elles détectent les hallucinations au niveau des \textit{tokens} ou des \textit{spans}, c'est-à-dire de segments continus de texte correspondant à une entité ou à une expression d'intérêt, plutôt qu'à l'échelle de la phrase entière. Zhou et al.~\cite{zhou2021detecting} entraînent ainsi un modèle à prédire, pour chaque \textit{token} généré, s'il est appuyé ou non par le texte source. Mishra et al.~\cite{mishra2024fava} proposent une taxonomie fine en six types d'hallucinations : les contradictions portant sur une entité, sur une relation ou sur une phrase entière, les entités inventées, les affirmations subjectives et les affirmations invérifiables. Leur analyse montre que les erreurs d'entité sont les plus fréquentes, représentant entre 42,7\,\% et 55,7\,\% des hallucinations selon le modèle. Ils entraînent ensuite FAVA, un modèle augmenté par récupération de documents qui repère les segments hallucinés, en indique le type, puis propose une correction. Dans une logique voisine, FActScore~\cite{min2023factscore} décompose un texte généré en faits atomiques, c'est-à-dire en affirmations élémentaires, et vérifie chacun d'eux de manière indépendante contre une source de connaissances. D'autres travaux exploitent les représentations internes des modèles de langage afin de repérer des comportements atypiques associés aux hallucinations. Azaria et Mitchell~\cite{azaria2023internal} montrent qu'un classifieur simple appliqué aux activations cachées d'un LLM permet de prédire si une affirmation produite est vraie ou fausse. Chen et al.~\cite{chen2024inside} proposent INSIDE, qui mesure la cohérence sémantique de plusieurs réponses directement dans l'espace des représentations internes du modèle, à l'aide d'un score fondé sur les valeurs propres de la matrice de covariance de leurs plongements (\textit{EigenScore}), et tronquent les activations extrêmes afin de mieux repérer les hallucinations produites avec une confiance excessive. L'idée est que les représentations produites lors de la génération diffèrent lorsque le modèle s'appuie sur des informations présentes dans le document ou, au contraire, lorsqu'il génère un contenu non fondé sur celui-ci. Ces dernières approches supposent toutefois un accès aux états internes du modèle, ce qui n'est pas possible lorsque celui-ci est utilisé en boîte noire à travers une API. Enfin, Kadavath et al.~\cite{kadavath2022know} montrent que les LLMs sont en partie capables d'estimer eux-mêmes la probabilité que leur réponse soit correcte, mais cette auto-évaluation reste imparfaite et dépend de la tâche et de la formulation de la question.

Un travail plus récent, LATTE~\cite{zhang2026latte}, propose une nouvelle architecture visant à améliorer la stabilité de l'extraction conjointe des entités et des relations, notamment en présence de triplets qui se chevauchent, c'est-à-dire de plusieurs relations partageant une même entité. Les auteurs obtiennent des scores F1 de 93,8\,\% et 95,1\,\% sur les corpus NYT et WebNLG. Cette étude montre que les erreurs liées à l'identification des entités et des relations restent un défi important. Toutefois, leur approche repose sur une modification de l'architecture du modèle pendant son entraînement, ce qui n'est pas applicable à un modèle multimodal propriétaire utilisé en boîte noire.

Ces travaux mettent en évidence un manque : à notre connaissance, il n'existe pas encore de caractérisation systématique des hallucinations propres aux schémas d'extraction contraints, où l'erreur peut porter indépendamment sur la valeur d'une entité, son type, ou la cohérence structurelle du triplet entité--relation--entité qu'elle compose. La taxonomie de Mishra et al.~\cite{mishra2024fava}, qui distingue erreurs d'entité et erreurs de relation, s'en approche, mais elle porte sur du texte libre et ne tient pas compte des contraintes d'un schéma.

\subsection{Hallucination dans les MLLMs appliqués aux documents scannés}

Ce sous-domaine constitue le croisement le plus proche des configurations d'extraction directe par MLLM, sans passage par un OCR intermédiaire. Les premiers travaux sur l'hallucination des modèles vision-langage ont porté sur les images naturelles : Rohrbach et al.~\cite{rohrbach2018chair} ont introduit la métrique CHAIR, qui mesure la proportion d'objets mentionnés dans une légende mais absents de l'image, et Li et al.~\cite{li2023pope} ont proposé POPE, qui évalue l'hallucination d'objets en interrogeant le modèle par des questions fermées (\og y a-t-il un \ldots{} dans l'image ? \fg{}). Ces protocoles ne s'appliquent toutefois pas directement aux documents textuels, où l'unité d'information n'est pas un objet mais une valeur ou une mention écrite.

He et al.~\cite{he2025seeing} proposent KIE-HVQA, présenté comme le premier benchmark dédié à l'hallucination OCR dans la compréhension de documents dégradés. Il comprend 2\,000 exemples d'entraînement et 400 exemples de test, répartis entre reçus, documents imprimés et cartes d'identité, auxquels sont appliquées des dégradations simulant des conditions réelles : flou de mouvement, faible contraste, occultation ou masquage d'une partie du texte. Chaque caractère est annoté selon sa lisibilité. Les auteurs montrent que, face à ces dégradations, les MLLMs reconnaissent rarement leur propre incertitude. Au lieu de signaler qu'un caractère est illisible, ils s'appuient sur des biais linguistiques, c'est-à-dire sur les connaissances apprises par le modèle et sur les formulations les plus probables, et produisent des valeurs absentes de l'image. Pour réduire ce phénomène, ils entraînent Qwen2.5-VL-7B par un affinage supervisé suivi d'un apprentissage par renforcement (GRPO), avec une récompense qui distingue les caractères lisibles, partiellement masqués et illisibles. Le modèle apprend ainsi à signaler explicitement les zones illisibles au lieu de les compléter. Le modèle obtenu dépasse GPT-4o d'environ 28 points de précision sans hallucination sur KIE-HVQA, tout en conservant ses performances OCR générales. Ce travail est proche du contexte des comptes rendus scannés : les auteurs citent explicitement le cas d'ordonnances délavées, dont l'unité de dosage est lisible mais pas la fréquence d'administration, et montrent que, pour un document dégradé, la bonne réponse est parfois de signaler l'incertitude plutôt que d'extraire une valeur. Il repose toutefois sur un réentraînement du modèle, ce qui n'est pas applicable à un modèle propriétaire utilisé en boîte noire.

Sur le plan mécanistique, Baek et al.~\cite{baek2025ocrheads} identifient au sein des modèles vision-langage des têtes d'attention spécialisées dans la lecture du texte présent dans l'image, qu'ils appellent têtes OCR. Ces têtes se distinguent des têtes de récupération (\textit{retrieval heads}) qui copient l'information dans un contexte textuel : elles se recouvrent peu avec elles, sont beaucoup moins rares, et leur activation est statique, c'est-à-dire directement liée à leur rôle de lecture, alors que les têtes de récupération sont activées de manière plus variable. Masquer une vingtaine de têtes OCR dans Qwen2-VL-7B fait chuter le F1 sur DocVQA de 76,3 à 59,7, ce qui confirme leur rôle central dans la lecture documentaire. Ces résultats offrent une piste de compréhension du comportement interne du modèle en complément des approches purement comportementales. Ils supposent cependant un accès aux poids et aux états internes du modèle, et les auteurs ne les évaluent pas directement pour la réduction des hallucinations.

Le travail le plus proche de la configuration visée par une extraction clinique directe est celui de RAPTOR+~\cite{raptorplus2026}, qui cible le traitement de formulaires scannés de signalement urgent pour suspicion de cancer colorectal, souvent de mauvaise qualité (tampons administratifs, ratures, mentions surchargées, modèles hétérogènes d'un cabinet adresseur à l'autre). Chaque champ extrait y est associé à une boîte englobante qui localise sa preuve dans l'image. Les auteurs y proposent un cadre d'évaluation quantifiant conjointement la précision de l'extraction structurée et la qualité de localisation de la preuve visuelle, et caractérisent deux modes d'échec cliniquement critiques : les extractions non ancrées, où le modèle produit une valeur plausible absente du document, et les prédictions invérifiables, où la valeur peut être correcte mais sans que le système puisse indiquer où se trouve la preuve. Sur 223 formulaires, Gemini 2.5 Flash, utilisé sans affinage, lit correctement 92,6\,\% des valeurs, mais seulement 1,2\,\% de ses extractions sont à la fois correctes et associées à une boîte de preuve bien placée (critère de sécurité strict). Après affinage, Qwen3-VL-8B atteint respectivement 96,1\,\% et 60,6\,\%. Ces résultats montrent qu'une bonne exactitude de lecture ne garantit pas un ancrage fiable, et que cet ancrage nécessite un affinage spécifique à la tâche.

Le principe de l'ancrage visuel, qui consiste à appuyer les réponses du modèle sur des éléments réellement présents dans l'image, a été mis en œuvre par Woodpecker~\cite{woodpecker2023} pour corriger les hallucinations d'un MLLM. Cette approche utilise des boîtes englobantes (\textit{bounding boxes}), c'est-à-dire des rectangles qui délimitent la position des objets ou des régions d'intérêt dans une image, comme preuves visuelles. Elle s'appuie sur un pipeline en cinq étapes : l'extraction de concepts clés, la formulation de questions, la validation des connaissances visuelles par des modèles experts, la génération d'affirmations visuelles, puis la correction des hallucinations. Woodpecker présente l'intérêt de ne nécessiter aucun réentraînement : la correction intervient après la génération, ce qui la rend applicable à un modèle utilisé en boîte noire. Dans le domaine médical, MedSIGHT~\cite{chang2026medsight} poursuit un objectif voisin par une autre voie : plutôt que de corriger la réponse après génération, il entraîne un modèle vision-langage médical à ancrer sa compréhension dans des régions de l'image, en produisant des masques de segmentation dont peuvent être dérivées des boîtes englobantes. Son application reste toutefois limitée à l'imagerie médicale (scanner, IRM, radiographie, ainsi que des images de pathologie pour les questions-réponses) et ne concerne pas les documents textuels scannés.

L'analyse de ces travaux met en évidence une limite : aucun ne combine à la fois un document clinique textuel scanné et dégradé, une extraction structurée des entités et des relations, ainsi qu'un mécanisme d'ancrage dans la source. Parmi les approches existantes, RAPTOR+~\cite{raptorplus2026} est la plus proche de cette configuration. Toutefois, elle ne produit pas de graphe de connaissances et n'intègre ni validation symbolique ni validation ontologique.

\subsection{Hallucination dans l'extraction d'information clinique à partir de texte propre}

La section~\ref{sec:extraction_er} a montré que les LLMs peuvent extraire des informations cliniques sans entraînement spécifique, à condition de guider le format de leurs réponses et de les post-traiter~\cite{agrawal2022fewshot}.

Ce contrôle est d'autant plus nécessaire que les LLMs hallucinent aussi dans le domaine médical. Pal et al.~\cite{pal2023medhalt} le mettent en évidence avec Med-HALT, un benchmark qui évalue les hallucinations des LLMs selon deux familles de tests : des tests de raisonnement, construits à partir de questions d'examens de médecine de plusieurs pays (Espagne, Inde, États-Unis, Taïwan), et des tests de mémoire, qui vérifient la capacité du modèle à restituer correctement des informations bibliographiques issues de PubMed. Wang et al.~\cite{wang2024feature} identifient les \textit{feature hallucinations}, c'est-à-dire la génération de caractéristiques cliniques incorrectes ou inexistantes, comme l'un des principaux obstacles à l'utilisation des LLMs pour extraire des données cliniques. Pour les limiter, ils proposent un pipeline modulaire qui décompose l'extraction en étapes successives (extraction des concepts, agrégation, génération de questions, extraction du passage pertinent, puis question-réponse), et l'évaluent sur les dossiers d'admission de 25\,709 grossesses. Les auteurs relèvent notamment que, sans extraction préalable du passage pertinent, le modèle confond les antécédents avec l'histoire de la maladie actuelle, ou les antécédents du patient avec ceux de sa famille. Avec ce pipeline, Qwen-14B-Chat et Baichuan2-13B-Chat atteignent respectivement 95,52\,\% et 95,86\,\% d'exactitude. De leur côté, Yang et al.~\cite{yang2025clines} soulignent, pour motiver leur système CLINES, que l'association des entités extraites par les LLMs aux terminologies médicales UMLS (\textit{Unified Medical Language System}) et SNOMED CT (\textit{Systematized Nomenclature of Medicine -- Clinical Terms}), deux référentiels utilisés pour normaliser les concepts cliniques, reste souvent imprécise et peut produire des codes hallucinés, et que les limites de longueur du contexte compliquent le raisonnement sur les notes cliniques volumineuses. CLINES y répond par un pipeline modulaire : découpage sémantique des notes longues, extraction par un LLM de raisonnement, attribution d'attributs (statut d'assertion, sujet de l'information, valeurs et unités), normalisation explicite vers l'UMLS et résolution des dates. Sur des notes de MIMIC-III, du consortium 4CE et des comptes rendus d'oncologie, il dépasse toutes les méthodes de référence, avec un gain de 0,21 à 0,38 de F1 sur un LLM interrogé par une seule invite. Les auteurs reconnaissent toutefois que le risque d'hallucination persiste, en particulier pour les abréviations rares : \og OD 2 gtt qd \fg{}, une posologie de collyre, est par exemple lu comme un surdosage (\textit{overdose}).

Dans une tâche proche, celle du résumé de dialogues entre médecin et patient, BN et al.~\cite{bn2025fact} montrent que les détecteurs d'hallucinations généralistes se comportent mal sur des textes cliniques, et que les résultats obtenus sur des hallucinations provoquées artificiellement ne prédisent pas ceux obtenus sur des hallucinations naturelles. Ils développent en revanche des détecteurs fondés sur l'atomisation des faits, qui consiste à décomposer un texte en faits élémentaires puis à vérifier chacun d'eux contre la transcription : ces détecteurs comptent les hallucinations de façon explicable et, bien que mis au point sur des hallucinations provoquées, se généralisent bien aux hallucinations naturelles. Enfin, Zeba et al.~\cite{zeba2026mitigating} proposent un module de vérification factuelle qui ne fait pas appel à un LLM : les résumés produits et le dossier du patient sont découpés en propositions élémentaires, comparées une à une au moyen de règles portant sur la négation, l'implication, la cohérence temporelle et les valeurs numériques. Ce travail montre qu'un mécanisme déterministe, c'est-à-dire reposant sur des règles produisant toujours le même résultat pour une même entrée, combinant plusieurs vérifications de cohérence logique, peut constituer une alternative intéressante.

\subsection{Auto-correction et correction après génération}

Une autre famille de travaux cherche non plus seulement à détecter les hallucinations, mais à les corriger après la génération. Madaan et al.~\cite{madaan2023selfrefine} proposent Self-Refine, dans lequel le modèle critique sa propre réponse puis la réécrit de manière itérative. Dhuliawala et al.~\cite{dhuliawala2024cove} proposent Chain-of-Verification : le modèle génère d'abord une réponse, puis formule des questions de vérification, y répond indépendamment et révise sa réponse initiale en conséquence.

La portée de ces approches est cependant discutée. Huang et al.~\cite{huang2024cannot} montrent qu'un LLM ne parvient pas à améliorer son raisonnement par auto-correction lorsqu'il ne dispose d'aucun retour extérieur, et que ses performances peuvent même se dégrader. La revue critique de Kamoi et al.~\cite{kamoi2024selfcorrection} conclut qu'aucun travail antérieur ne démontre une auto-correction réussie à partir d'un retour produit par le LLM lui-même sur des tâches générales, sauf pour des tâches particulièrement favorables, et que l'auto-correction n'est fiable que lorsque le modèle reçoit un retour provenant d'une source extérieure fiable, par exemple un exécuteur de code ou une recherche sur le web, ou après un affinage à grande échelle. Ces résultats sont importants pour l'extraction d'information : demander au modèle de vérifier lui-même ses extractions ne garantit pas leur fidélité au document, et la vérification gagne à s'appuyer sur des signaux extérieurs au modèle, tels que le texte source ou les contraintes d'un schéma.

\subsection{Ontologies, contraintes et réparation de graphes comme mécanisme de validation}

La vérification de la cohérence structurelle d'un graphe est un domaine bien établi en représentation des connaissances. La recommandation SHACL (\textit{Shapes Constraint Language}) du W3C~\cite{shacl2017} permet de déclarer les contraintes qu'un graphe RDF doit respecter, telles que les types autorisés, les propriétés attendues, les cardinalités ou le domaine et l'image d'une relation, et produit un rapport listant chaque violation détectée. Ce type de validation détecte les incohérences, mais ne les répare pas. En outre, un graphe peut respecter toutes les contraintes déclarées tout en contenant une information absente du document : la validité structurelle ne garantit pas la fidélité à la source.

La réparation des erreurs a été étudiée principalement dans le domaine des bases de données et des graphes de connaissances. Paulheim~\cite{paulheim2017refinement} recense les méthodes de raffinement de graphes de connaissances, qu'il répartit en deux familles : la complétion, qui ajoute les faits ou les types manquants, et la détection d'erreurs, qui repère les assertions erronées. La correction proprement dite des erreurs y reste peu traitée. Il souligne également la difficulté d'évaluer ces méthodes en l'absence de vérité de terrain complète. En bases de données, la réparation consiste à rechercher la modification minimale qui fait disparaître les violations de contraintes. Chu et al.~\cite{chu2013holistic} traitent l'ensemble des violations simultanément, puisqu'une même valeur erronée peut en provoquer plusieurs, Fan et al.~\cite{fan2016gfd} définissent des dépendances fonctionnelles adaptées aux graphes, et HoloClean~\cite{rekatsinas2017holoclean} combine contraintes, statistiques et inférence probabiliste pour proposer des réparations. Ces méthodes ont toutefois été conçues pour des données déjà établies, où les erreurs sont rares et isolées, et elles ne s'appuient pas sur un document source pour justifier une correction.

Plusieurs travaux récents mobilisent par ailleurs les ontologies biomédicales comme mécanisme de validation indépendant du modèle générateur. Par exemple, Das et al.~\cite{das2026clinical} construisent un graphe de connaissances clinique à partir de comptes rendus d'oncologie à l'aide d'un pipeline combinant une extraction guidée par \textit{prompts}, un score d'incertitude fondé sur l'entropie, la génération d'un schéma RDF/OWL aligné sur des ontologies médicales (SNOMED CT, LOINC, RxNorm) et une validation par consensus multi-LLM, où plusieurs modèles comparent leurs résultats afin d'obtenir une décision plus fiable. L'évaluation est toutefois menée sans annotation de référence. Ali et al.~\cite{ali2026ontology} proposent un cadre GraphRAG : ils construisent une ontologie RDF/OWL et un graphe de connaissances à partir de données cliniques et hospitalières issues de plusieurs établissements égyptiens, puis les intègrent à un LLM pour ancrer ses réponses dans une tâche de question-réponse clinique. Sur 60 questions, l'exactitude passe de 37\,\% (ChatGPT-4) et 52\,\% (DeepSeek-R1) à 98\,\%, et le taux d'hallucination de 63\,\% et 48\,\% à 1,7\,\%. Ces résultats restent préliminaires, compte tenu de la taille réduite de l'évaluation. Ce travail porte en outre sur la génération de réponses, et non sur la validation d'un graphe extrait de documents.

En dehors du domaine médical, CyberBOT~\cite{zhao2025cyberbot}, un assistant pédagogique en cybersécurité, ajoute une étape de vérification après la génération : les réponses d'un système RAG sont confrontées à une ontologie du domaine construite par des experts, et celles jugées incohérentes avec ses concepts et ses relations sont rejetées. Cette vérification est toutefois elle-même réalisée par un LLM (Llama 3.3 70B), et non par un raisonneur formel, et les auteurs ne mesurent pas son effet propre. Ils observent d'ailleurs que la fidélité est le plus faible de leurs indicateurs (0,788 en moyenne), le modèle recourant à ses propres connaissances lorsque la base documentaire ne contient pas de passage pertinent. À l'inverse, Sadowski et Chudziak~\cite{sadowski2026structured} confient la vérification à des outils symboliques : le LLM décompose le raisonnement en étapes (identification des entités, extraction d'assertions dans une ontologie OWL), puis des règles SWRL appliquées par un raisonneur logique, c'est-à-dire un outil capable de vérifier automatiquement la cohérence des connaissances, produisent la décision. Sur trois domaines, dont l'éligibilité aux essais cliniques, ils obtiennent des gains statistiquement significatifs par rapport au \textit{prompting} à partir de quelques exemples (\textit{few-shot}), et une étude d'ablation confirme l'apport propre de la vérification symbolique.

Mandarapu et Kunkunuru~\cite{mandarapu2026knowledge} mènent une étude contrôlée sur l'ancrage par graphe de connaissances, qui consiste à interroger un graphe de connaissances pour fournir au modèle les faits nécessaires à sa réponse, dans une tâche de question-réponse clinique (MedQA). Sur le graphe biomédical public PrimeKG, ni une récupération naïve de triplets ni un agent générant ses propres requêtes sur le graphe n'améliorent significativement les performances, quel que soit le modèle de base (écarts compris entre $-3{,}4$ et $+1{,}4$ points, dans les intervalles de confiance). À l'inverse, sur un graphe synthétique contenant des faits fictifs, le même pipeline fait passer la précision du hasard (environ 22\,\%) à près de 100\,\%. Les auteurs en concluent que l'ancrage n'est utile que lorsque le fait décisif se situe hors des données d'entraînement du modèle et ne peut pas être déduit de sa forme, comme le nom d'un médicament qui renseigne souvent sur sa classe thérapeutique. Les faits d'un graphe public étant déjà connus des modèles, l'ancrage est redondant ; il devient décisif pour les données privées et institutionnelles, qu'aucun modèle n'a pu voir. Les auteurs rappellent par ailleurs qu'une étude antérieure, parue dans \textit{Nature Medicine}, observe que la récupération d'informations peut même dégrader les performances d'un modèle déjà performant.

\subsection{Hybrid-Code v2 : génération neuronale et validation symbolique indépendante}

Un travail neuro-symbolique récent, Hybrid-Code v2~\cite{hybrid2512v2}, combine un modèle neuronal avec un mécanisme de validation symbolique indépendant. Ce système génère d'abord des candidats de codes ICD-10, une classification internationale des maladies utilisée pour coder les diagnostics médicaux, puis les vérifie à l'aide d'une base de connaissances symbolique. La validation repose sur plusieurs niveaux de contrôle, notamment le format du code, son ancrage dans le document, la détection des négations, la cohérence temporelle et des règles d'exclusion. Les auteurs montrent que cette approche élimine complètement les erreurs de format (hallucinations de type I) et réduit les erreurs cliniquement incorrectes mais syntaxiquement valides (hallucinations de type II), sans toutefois les supprimer totalement, l'auteur précisant que la correction des erreurs de type II n'est pas formellement garantie. Évalué sur MIMIC-III, le système atteint une couverture de 85\,\% et une précision de 92\,\%, sans aucune hallucination de type I, alors que les modèles neuronaux de référence en produisent entre 6 et 18\,\%.

Ce travail montre qu'il est possible de combiner une génération neuronale avec une validation symbolique indépendante. Toutefois, il s'intéresse à une tâche de classification fermée, c'est-à-dire au choix d'un code parmi une nomenclature existante de plus de 70\,000 codes (ICD-10-CM), à partir d'un texte clinique déjà propre. Contrairement aux tâches d'extraction de connaissances, il ne nécessite ni l'identification des entités et des relations ni la construction d'un graphe de connaissances.

\subsection{Architectures multi-agents pour la détection d'hallucination clinique}

L'idée de faire collaborer plusieurs modèles pour améliorer la fiabilité des réponses a d'abord été étudiée hors du domaine médical. Du et al.~\cite{du2024debate} montrent que faire débattre plusieurs instances d'un LLM, chacune critiquant les réponses des autres sur plusieurs tours, améliore la factualité et le raisonnement par rapport à un modèle seul. Hong et al.~\cite{hong2024metagpt} observent à l'inverse que l'enchaînement naïf de plusieurs LLMs peut provoquer des hallucinations en cascade. Leur cadre MetaGPT y répond en encodant des procédures opératoires standardisées (\textit{Standardized Operating Procedures}) dans les invites, en attribuant à chaque agent un rôle précis et en imposant des sorties intermédiaires structurées, plutôt qu'un dialogue libre entre agents.

Dans le domaine clinique, CuraView~\cite{ye2026curaview}, Dialectic-Med~\cite{lu2026dialectic} et BLUEmed~\cite{you2026bluemed} proposent des architectures multi-agents, dans lesquelles plusieurs agents aux rôles complémentaires (génération, recherche de preuves, contradiction, arbitrage) collaborent pour détecter ou atténuer les hallucinations. CuraView vérifie, phrase par phrase, la fidélité de comptes rendus de sortie générés par un LLM au dossier du patient, à l'aide d'un graphe de connaissances construit à partir de ce dossier (GraphRAG) et d'une gradation des niveaux de preuve. Dialectic-Med fait débattre un agent proposant, un agent opposant et un médiateur afin de limiter les hallucinations diagnostiques des MLLMs médicaux. BLUEmed combine recherche augmentée et débat entre agents experts pour détecter les erreurs terminologiques dans des notes cliniques. Ces approches portent toutefois sur la détection ou l'atténuation d'erreurs dans des textes générés ou des diagnostics, et non sur l'extraction structurée d'entités et de relations, et elles n'organisent pas leurs décisions selon une typologie des erreurs d'extraction. De plus, lorsque les agents vérificateurs sont eux-mêmes des modèles de langage, la vérification hérite du risque qu'elle cherche à contrôler : un agent juge peut lui aussi halluciner. Ce constat rejoint les conclusions de Kamoi et al.~\cite{kamoi2024selfcorrection} sur la nécessité d'un retour extérieur fiable.

\subsection{Vérification sans modèle juge : cohérence, contexte clinique et décision sélective}

Face à ce risque de circularité, une autre famille de travaux, plus ancienne pour certains, repose sur des signaux explicites et reproductibles : la cohérence entre plusieurs générations, la détection de règles linguistiques dans le texte clinique, la combinaison pondérée d'heuristiques et la possibilité pour un système de s'abstenir. Ces travaux ne ciblent pas tous directement l'hallucination des MLLMs, mais ils fournissent des briques méthodologiques transposables à l'extraction structurée.

\paragraph{Cohérence entre générations.}
Wang et al.~\cite{wang2023selfconsistency} proposent la \textit{self-consistency} : au lieu de conserver une seule réponse, le modèle génère plusieurs chemins de raisonnement, puis la réponse la plus fréquente est retenue par vote majoritaire. Les auteurs montrent que cette stratégie améliore les performances sur des tâches de raisonnement arithmétique et de sens commun. Manakul et al.~\cite{manakul2023selfcheckgpt} exploitent une intuition voisine pour la détection d'hallucinations avec SelfCheckGPT : si le modèle connaît réellement une information, plusieurs générations stochastiques tendent à être cohérentes entre elles, alors qu'une information hallucinée tend à varier d'une génération à l'autre. Farquhar et al.~\cite{farquhar2024semantic} affinent cette idée avec l'entropie sémantique, qui regroupe les générations ayant le même sens avant de mesurer leur dispersion, afin de ne pas confondre une simple variation de formulation avec une incertitude réelle. Ces approches fonctionnent sans base de connaissances externe et, pour SelfCheckGPT, en boîte noire, ce qui les rend compatibles avec un modèle propriétaire accessible uniquement par API, au prix de plusieurs appels au modèle. Toutefois, la cohérence ne garantit pas la fidélité au document : un modèle peut produire de manière répétée la même information plausible mais absente de la source. Dans un contexte d'extraction, le vote doit donc être combiné à une vérification d'ancrage dans le document.

\paragraph{Négation, hypothèse, temporalité et sujet de l'information.}
Chapman et al.~\cite{chapman2001negex} proposent NegEx, un algorithme fondé sur des expressions régulières qui détermine si un concept clinique est nié à partir de déclencheurs (par exemple \og no \fg{} ou \og without \fg{}) situés dans une fenêtre de mots autour de la mention, initialement sur des comptes rendus de sortie. Harkema et al.~\cite{harkema2009context} étendent ce principe avec ConText, qui détermine trois propriétés contextuelles d'une mention : la négation, le statut temporel (récent, historique ou hypothétique) et le sujet de l'information (\textit{experiencer}), c'est-à-dire le fait que le constat concerne le patient ou une autre personne. Ces algorithmes sont déterministes et interprétables. Leur intérêt pour l'étude des hallucinations tient au fait qu'ils permettent de détecter un type d'erreur invisible à une simple vérification de présence dans le texte : un modèle qui extrait \og fièvre \fg{} à partir de \og absence de fièvre \fg{} produit une entité littéralement présente dans le document, mais un fait clinique faux. Ces méthodes ont été conçues pour l'anglais ; leur application au français suppose l'adaptation des listes de déclencheurs, comme l'ont fait Deléger et Grouin~\cite{deleger2012negation}, qui ont construit une liste de déclencheurs en français à partir de notes cliniques et obtenu une F-mesure de 0,867, comparable à celle de NegEx sur l'anglais (0,839).

\paragraph{Combinaison pondérée d'heuristiques.}
Ratner et al.~\cite{ratner2017snorkel} proposent Snorkel, un système de supervision faible dans lequel des experts écrivent des fonctions d'étiquetage (\textit{labeling functions}), c'est-à-dire des règles heuristiques bruitées, dont les sorties sont combinées par un modèle qui estime la fiabilité de chacune afin de produire des étiquettes d'entraînement. L'intérêt de ce travail, au-delà de la création de données d'entraînement, réside dans l'idée qu'un ensemble de règles individuellement imparfaites peut être agrégé en une décision plus fiable, à condition d'estimer le poids de chaque règle plutôt que de les traiter toutes de manière identique.

\paragraph{Décision sélective.}
L'idée de permettre à un système de refuser de décider est ancienne : Chow~\cite{chow1970reject} formalise dès 1970 le compromis entre le taux d'erreur et le taux de rejet d'un classifieur. Geifman et El-Yaniv~\cite{geifman2017selective} transposent ce principe aux réseaux de neurones profonds avec la classification sélective : le système peut s'abstenir lorsque sa confiance est insuffisante, et un seuil sur une fonction de confiance est choisi de façon à garantir, avec une forte probabilité, un niveau de risque fixé à l'avance, tout en conservant la plus grande couverture possible. Ce compromis entre risque et couverture est bien adapté au domaine médical, où il est souvent préférable de ne pas produire une information plutôt que de produire une information fausse. Transposé à l'extraction, il correspond à un compromis explicite entre précision et rappel.

\paragraph{Bilan.}
Ces quatre familles de travaux couvrent chacune une partie du problème : la cohérence entre générations cible l'instabilité du modèle, les algorithmes de type NegEx et ConText ciblent les erreurs de contexte et de temporalité, la supervision faible fournit un cadre pour combiner des indices hétérogènes, et la décision sélective permet de contrôler explicitement le niveau de risque.

\subsection{Benchmarks et métriques d'évaluation}

Plusieurs benchmarks, c'est-à-dire des jeux de données et des protocoles de référence utilisés pour évaluer les modèles, sont aujourd'hui consacrés à l'évaluation des hallucinations en contexte clinique. Med-HALT~\cite{pal2023medhalt} évalue les LLMs sur des tests de raisonnement issus d'examens de médecine de plusieurs pays et sur des tests de mémoire fondés sur PubMed. ClinHallu~\cite{yang2026clinhallu}, destiné aux MLLMs médicaux, identifie l'étape du raisonnement où apparaît l'hallucination (perception, connaissance ou intégration) au lieu d'évaluer uniquement la réponse finale. Pour les modèles multimodaux, CHAIR~\cite{rohrbach2018chair} et POPE~\cite{li2023pope} mesurent l'hallucination d'objets dans des images naturelles, et KIE-HVQA~\cite{he2025seeing} l'hallucination de lecture sur des documents dégradés. Toutefois, aucun n'est spécifiquement conçu pour l'extraction d'entités et de relations à partir de documents cliniques scannés présentant de fortes contraintes structurelles.

Concernant les métriques d'évaluation, la plupart des travaux d'extraction se limitent à la précision, au rappel et au F1-score. Or, la précision mesure indirectement le taux d'hallucination, puisque toute information extraite absente de la référence est comptée comme un faux positif, mais elle ne distingue pas une information inventée d'une information correcte mal typée ou mal délimitée. Des métriques plus fines ont été proposées : FActScore~\cite{min2023factscore} mesure la proportion de faits atomiques appuyés par une source, et Ghanem et Cruz~\cite{ghanem2024enhancing} proposent un cadre reposant sur une mesure de similarité entre graphes basée sur BERTScore, une métrique qui compare la similarité sémantique entre deux textes à l'aide de représentations produites par BERT. Les auteurs recommandent un seuil de 95\,\% pour considérer que deux graphes correspondent, dans le cadre de la réduction des hallucinations et des omissions lors de la construction de graphes de connaissances à partir de textes génériques. Leurs expériences, menées avec le modèle Mistral, montrent qu'un affinage réduit les hallucinations et les omissions exactes, mais dégrade la généralisation à un autre jeu de données. Bien que cette étude ne porte pas sur le domaine clinique, elle constitue une référence importante pour l'évaluation de la fidélité structurelle, au-delà des métriques classiques telles que la précision et le rappel.

\subsection{Analyse critique et synthèse}

L'analyse de la littérature fait apparaître une couverture inégale des différentes dimensions du problème. L'ancrage visuel sur documents dégradés est partiellement couvert, principalement par RAPTOR+~\cite{raptorplus2026}, mais hors du cadre d'une extraction structurée vers un graphe de connaissances. La validation ontologique est bien établie sur texte propre, et la réparation sous contraintes est bien étudiée en bases de données, mais sans exiger que la correction soit justifiée par un document source. Les méthodes d'auto-correction et les architectures multi-agents introduisent des mécanismes de vérification, mais sans typologie fine des erreurs et en s'appuyant le plus souvent sur des modèles de langage comme juges, dont la fiabilité dépend de l'existence d'un retour extérieur~\cite{kamoi2024selfcorrection}. Hybrid-Code v2~\cite{hybrid2512v2} atteint une garantie formelle sur une catégorie d'erreur, mais sur une tâche de classification fermée qui ne pose jamais la question de la frontière entre entité et relation. Enfin, les méthodes de vérification sans modèle juge restent étudiées isolément, rarement au sein d'un pipeline de construction de graphe de connaissances à partir d'un MLLM.

Cette cartographie met ainsi en évidence un espace de recherche encore ouvert à l'intersection de trois éléments rarement combinés dans la littérature actuelle : le traitement direct de documents cliniques scannés, l'extraction structurée d'entités et de relations vers un graphe de connaissances, et une validation combinant ancrage dans la source, cohérence symbolique et vérification explicite de la stabilité de la frontière entité--relation. LATTE~\cite{zhang2026latte} aborde bien cette instabilité, mais au niveau de l'architecture du modèle et non en validation \textit{post-hoc}.


\subsection{Synthèse des hallucinations identifiées dans la littérature}

Le Tableau~\ref{tab:hallucinations} synthétise les principaux types d'hallucinations identifiés dans cette section, les travaux qui les étudient et les approches proposées pour les réduire. Il montre que les hallucinations ne se limitent pas à l'invention d'informations : elles peuvent affecter la perception visuelle, les entités, les relations, le contexte clinique de la mention, le raisonnement ou la cohérence structurelle du graphe.


{\footnotesize
\renewcommand{\arraystretch}{1.35}
\begin{longtable}{|>{\raggedright\arraybackslash}p{2.6cm}|>{\raggedright\arraybackslash}p{4.2cm}|>{\raggedright\arraybackslash}p{3.4cm}|>{\raggedright\arraybackslash}p{4.0cm}|}
\caption{Synthèse des principaux types d'hallucinations identifiés dans la littérature et des approches proposées pour les réduire.}
\label{tab:hallucinations}\\
\hline
\textbf{Type d'hallucination} & \textbf{Description} & \textbf{Travaux} & \textbf{Approches proposées} \\
\hline
\endfirsthead
\multicolumn{4}{l}{\textit{Tableau~\ref{tab:hallucinations} (suite)}}\\
\hline
\textbf{Type d'hallucination} & \textbf{Description} & \textbf{Travaux} & \textbf{Approches proposées} \\
\hline
\endhead
\hline
\multicolumn{4}{r}{\textit{Suite à la page suivante}}\\
\endfoot
\hline
\endlastfoot

Hallucination visuelle et de lecture &
Le modèle génère des informations absentes du document ou interprète incorrectement les éléments visuels, en particulier lorsque la qualité de l'image est dégradée. &
CHAIR \cite{rohrbach2018chair}, POPE \cite{li2023pope}, KIE-HVQA \cite{he2025seeing}, Woodpecker \cite{woodpecker2023}, MedSIGHT \cite{chang2026medsight}, RAPTOR+ \cite{raptorplus2026} &
Ancrage visuel, boîtes englobantes (\textit{bounding boxes}), segmentation, abstention sur les caractères illisibles et vérification après génération.\\
\hline

Hallucination d'entité &
Erreur sur la valeur d'une entité ou sur son type. &
LATTE \cite{zhang2026latte}, Wang et al. \cite{wang2024feature}, FAVA \cite{mishra2024fava} &
Amélioration de l'extraction, stabilisation des frontières de décision, détection au niveau des segments, validation des entités extraites.\\
\hline

Hallucination de relation &
Relation incorrecte entre deux entités ou mauvaise association entre les concepts extraits. &
LATTE \cite{zhang2026latte}, FAVA \cite{mishra2024fava} &
Extraction conjointe entité--relation, amélioration de la stabilité des frontières et détection des relations erronées.\\
\hline

Hallucination contextuelle &
Fait extrait alors qu'il est nié, hypothétique ou qu'il ne concerne pas le patient. &
NegEx \cite{chapman2001negex}, ConText \cite{harkema2009context}, Deléger et Grouin \cite{deleger2012negation}, Hybrid-Code v2 \cite{hybrid2512v2} &
Détection de déclencheurs de négation et d'hypothèse, identification du sujet de l'information.\\
\hline

Hallucination temporelle &
Information historique (antécédent) interprétée comme un fait actuel. &
ConText \cite{harkema2009context}, Hybrid-Code v2 \cite{hybrid2512v2} &
Détermination du statut temporel de la mention, vérification de la cohérence temporelle.\\
\hline

Hallucination ontologique &
Association erronée entre une entité extraite et un concept médical normalisé (UMLS, SNOMED CT ou ICD-10). &
CLINES \cite{yang2025clines}, Das et al. \cite{das2026clinical}, Hybrid-Code v2 \cite{hybrid2512v2} &
Validation ontologique, bases de connaissances symboliques, consensus multi-LLM et vérification des codes.\\
\hline

Hallucination de raisonnement &
Erreur apparaissant lors du raisonnement clinique ou de l'intégration des informations. &
Med-HALT \cite{pal2023medhalt}, ClinHallu \cite{yang2026clinhallu}, Dialectic-Med \cite{lu2026dialectic} &
Débat multi-agents et évaluation étape par étape du raisonnement.\\
\hline

Hallucination factuelle &
Production de faits médicaux inexacts ou omission d'informations pertinentes malgré un texte cohérent. &
BN et al. \cite{bn2025fact}, Zeba et al. \cite{zeba2026mitigating}, Wang et al. \cite{wang2024feature}, FActScore \cite{min2023factscore}, SelfCheckGPT \cite{manakul2023selfcheckgpt}, entropie sémantique \cite{farquhar2024semantic} &
Atomisation des faits, extraction modulaire, vérification factuelle déterministe et cohérence entre générations.\\
\hline

Hallucination structurelle &
Incohérences dans le graphe de connaissances, relations incompatibles ou non-respect des contraintes du schéma ou de l'ontologie. &
SHACL \cite{shacl2017}, Paulheim \cite{paulheim2017refinement}, HoloClean \cite{rekatsinas2017holoclean}, CyberBOT \cite{zhao2025cyberbot}, Sadowski et al. \cite{sadowski2026structured}, Ghanem et Cruz \cite{ghanem2024enhancing}, Hybrid-Code v2 \cite{hybrid2512v2} &
Validation par contraintes, réparation minimale, vérification ontologique après génération, raisonneurs logiques et mesure de similarité entre graphes.\\
\end{longtable}
}

\section{Conclusion du chapitre}

Ce chapitre a montré que les briques nécessaires à une extraction clinique fiable existent séparément dans la littérature : extraction multimodale, extraction d'entités et de relations par des modèles supervisés ou génératifs, annotation guidée par des règles explicites, ancrage dans la source, validation par contraintes et réparation de graphes, détection du contexte clinique, vote entre générations et décision sélective. Il a également montré que l'auto-correction par le modèle lui-même ou par d'autres modèles juges reste fragile en l'absence d'un retour extérieur fiable. En revanche, ces briques ont rarement été assemblées dans un même pipeline appliqué à des documents cliniques scannés et orienté vers la construction d'un graphe de connaissances. Le chapitre suivant présente l'approche proposée pour répondre à cette limite.
